\documentclass[../../../include/open-logic-section]{subfiles}

\begin{document}

\olfileid{sth}{z}{pairs}
\olsection{زوج‌ها}

اصل موضوع بعدی که باید بررسی کنیم چنین است:

\begin{axiom}[زوج‌ها]
برای هر دو مجموعهٔ~$a,b$، مجموعهٔ~$\{a,b\}$ وجود دارد.
\[
	\forall a \forall b \exists P \forall x (x \in P \liff (x = a \lor x = b))
\]
\end{axiom}

در ادامه این اصل موضوع را با استفاده از تصور تکراری توجیه می‌کنیم. فرض
کنید~$a$ در مرحلهٔ~$S$ و~$b$ در مرحلهٔ~$T$ در دسترس باشد. هرکدام از
مراحل~$S$ و~$T$ را که دیرتر می‌آید~$M$ بنامید. چون~$a$ و~$b$ هر دو در
مرحلهٔ~$M$ در دسترس‌اند، گردایهٔ~$\{a,b\}$ در هر مرحلهٔ پس از~$M$ گردایه‌ای
ممکن و در دسترس است.

اما لحظه‌ای درنگ کنید!{} چرا فرض کنیم اصلاً مرحله‌ای پس از~$M$ وجود دارد؟
اگر چنین مرحله‌ای نباشد، توجیه ما شکست می‌خورد. پس برای توجیه اصل زوج‌ها
باید اصل دیگری به روایتی بیفزاییم که در~\olref[sth][z][story]{sec} بیان
کردیم، یعنی:
\begin{enumerate}
	\item[] \stagessucc. هیچ مرحلهٔ آخری وجود ندارد.
\end{enumerate}
آیا این اصل موجه است؟ هیچ چیز در روایت شوئنفیلد \emph{به‌صراحت} نمی‌گفت
که مرحلهٔ آخری وجود ندارد. بااین‌حال، حتی اگر این اصل (به‌بیان دقیق) افزوده‌ای
بر روایت ما باشد، با اندیشهٔ پایه‌ایِ تشکیل مجموعه‌ها در مراحل سازگار است.
در ادامه صرفاً آن را می‌پذیریم. و در نتیجه، اصل موضوع زوج‌ها را نیز
خواهیم پذیرفت.

با دراختیارداشتن این اصل موضوع تازه می‌توانیم وجود مجموعه‌های بسیار دیگری
را اثبات کنیم. برای نمونه:

\begin{prop}\ollabel{prop:pairsconsequences}
برای هر دو مجموعهٔ~$a$ و~$b$، مجموعه‌های زیر وجود دارند:
	\begin{enumerate}
		\item\ollabel{singleton} $\{a\}$
		\item\ollabel{binunion} $a \cup b$
		\item\ollabel{tuples} $\tuple{a, b}$
	\end{enumerate}
\end{prop}

\begin{proof}
\olref{singleton}. بنا بر اصل زوج‌ها، $\{a,a\}$ وجود دارد که بنا بر
گسترش‌مندی همان~$\{a\}$ است.

\olref{binunion}. بنا بر اصل زوج‌ها، $\{a,b\}$ وجود دارد. اکنون بنا بر
اصل اجتماع، $a \cup b = \bigcup\{a,b\}$ وجود دارد.

\olref{tuples}. بنا بر~\olref{singleton}، $\{a\}$ وجود دارد. بنا بر اصل
زوج‌ها، $\{a,b\}$ وجود دارد. اکنون با یک کاربرد دیگر اصل زوج‌ها،
$\{\{a\},\{a,b\}\} = \tuple{a,b}$ وجود دارد.
\end{proof}

\begin{prob}
نشان دهید برای هر سه مجموعهٔ~$a,b,c$، مجموعهٔ~$\{a,b,c\}$ وجود دارد.
\end{prob}

\begin{prob}
نشان دهید برای هر مجموعه‌های~$a_1,\ldots,a_n$، مجموعهٔ
$\{a_1,\ldots,a_n\}$ وجود دارد.
\end{prob}

%\begin{proof} با دو کاربرد اصل زوج‌ها، $\{\{a_1,a_2\},\{a_1,
%   a_3\}\}$ وجود دارد. بنا بر اجتماع و گسترش‌مندی، $\bigcup \{\{a_1,
%   a_2\}, \{a_1,a_3\}\} = \{a_1,a_2,a_3\}$ وجود دارد. این شگرد را
%   هرچند بار که لازم است تکرار کنید. \end{proof}

\end{document}
